% =====================================================================
%  Bioinformática Práctica — libro de texto para maestría
%  Juvenal Yosa, PhD · juvenal.yosa@gmail.com · Copiloto: Claude (Anthropic)
%  Compilar:  latexmk -lualatex -shell-escape main.tex   (ver Makefile)
% =====================================================================
\documentclass[11pt,twoside,openright]{book}
\usepackage{estilo/bioinfo}
\usepackage{estilo/libro}

% Bibliografías por capítulo (se cargan solo las que existen)
\ExplSyntaxOn
\int_step_inline:nnn {0} {18} {
  \tl_set:Ne \l_tmpa_tl { bib/cap \int_compare:nNnT {#1} < {10} {0} #1 .bib }
  \file_if_exist:VT \l_tmpa_tl { \exp_args:NV \addbibresource \l_tmpa_tl }
}
\ExplSyntaxOff
% Capítulos (se incluyen solo los que existen)
\newcommand{\capitulo}[1]{\IfFileExists{capitulos/cap#1.tex}{\include{capitulos/cap#1}}{}}

\hypersetup{pdftitle={Bioinformática Práctica},
  pdfauthor={Juvenal Yosa},
  pdfsubject={Bioinformática para maestría},
  pdfkeywords={bioinformática, genómica, alineamiento, filogenética, NGS, RNA-seq}}

\begin{document}
\frontmatter
\portada

% Página de créditos (reverso de la portada interior)
\thispagestyle{empty}
\vspace*{\fill}
{\sffamily\small\color{tinta2}\setlength{\parindent}{0pt}\setlength{\parskip}{6pt}
{\titlefont\Large\color{azulnoche} Bioinformática Práctica}\\
{\itshape De la secuencia al sistema: fundamentos, algoritmos y laboratorio computacional}

\medskip
\textbf{Autor:} Juvenal Yosa, PhD · \href{mailto:juvenal.yosa@gmail.com}{juvenal.yosa@gmail.com}\\
\textbf{Copiloto:} Claude (Anthropic), asistente de inteligencia artificial empleado en la
redacción, la programación de figuras y la verificación bibliográfica, bajo la dirección y
revisión del autor.

\textbf{Laboratorios:} cada capítulo tiene notebooks de Google Colab en
\href{https://github.com/juvenalyosa/bioinformatica-practica}{\ttfamily github.com/\allowbreak juvenalyosa/\allowbreak bioinformatica-practica}.

\textbf{Primera edición:} 2026.

\textbf{Licencia:} texto, figuras y código se distribuyen bajo la licencia MIT. Se permite
usar, copiar, modificar y redistribuir este material, con o sin fines comerciales, siempre que se
conserve el aviso de autoría y de licencia.

\textbf{Referencias:} todas las referencias bibliográficas fueron verificadas contra los registros
de Crossref, DataCite u OpenLibrary (DOI o ISBN) y siguen las normas APA (7.ª edición).

\textbf{Tipografía:} compuesto con Lua\LaTeX{} en Libertinus Serif, Source Sans Pro y Fira Mono.
Figuras vectoriales en TikZ y PGFPlots, generadas a partir de cálculos reproducibles.

\medskip
\begin{tikzpicture}
  \foreach \c [count=\i] in {nucA,nucC,nucG,nucT}{\fill[\c] (\i*0.5,0) rectangle ++(0.42,0.12);}
\end{tikzpicture}
\par}
\cleardoublepage

\chapter*{Prólogo}
\addcontentsline{toc}{chapter}{Prólogo}
\markboth{Prólogo}{Prólogo}

En 1977, secuenciar unos pocos miles de bases del bacteriófago $\varphi$X174 fue una hazaña que
llevó años. Hoy, un solo instrumento produce en un día más secuencias de las que existían en todas
las bases de datos del mundo a comienzos de este siglo. La biología se ha convertido, casi sin
pedir permiso, en una ciencia de la información. Quien quiera hacer preguntas sobre genes, células,
poblaciones o ecosistemas necesita hoy dos alfabetos: el de las moléculas y el de los algoritmos.

Este libro nació como acompañante de un curso de bioinformática práctica para estudiantes de
maestría. Cada capítulo corresponde a un módulo del curso, y cada sección a una lección con un
laboratorio ejecutable en Google Colab. Pero el libro está pensado para leerse por sí mismo: allí
donde el notebook muestra \emph{cómo} hacer algo, el texto se detiene a explicar \emph{por qué}
funciona, de dónde viene cada ecuación y qué supuestos la sostienen.

El recorrido tiene una lógica deliberada. Cada tema comienza con una situación cotidiana que
ofrece un punto de apoyo intuitivo; sigue con las definiciones básicas y, paso a paso, sube hasta
la teoría formal que se encuentra en la literatura especializada. Las ecuaciones no aparecen como
adornos: cada una va acompañada de una tabla que explica todos sus símbolos, de un ejemplo
numérico resuelto y de un fragmento de código que la pone a trabajar. Las figuras se generaron a
partir de cálculos reales y reproducibles, y las referencias, que remiten siempre a las fuentes
originales, fueron verificadas una a una.

El libro está organizado en cinco partes. La primera prepara el laboratorio digital y el
vocabulario molecular. La segunda estudia la comparación de secuencias, desde el alineamiento
hasta la filogenética. La tercera sigue a las lecturas de secuenciación desde el instrumento hasta
las variantes. La cuarta aborda la genómica de poblaciones, la transcriptómica, la epigenómica y
la metagenómica. La quinta sube de escala hacia la estructura de las proteínas, las redes, el
aprendizaje automático y los flujos de trabajo reproducibles.

Nada de esto sustituye la práctica. La bioinformática se aprende como se aprende un instrumento:
leyendo la partitura, sí, pero sobre todo tocando. Por eso le invito a leer con el notebook abierto,
a cambiar los parámetros, a romper el código y a repararlo. Si al terminar un capítulo usted puede
explicar a un colega por qué un algoritmo funciona, y no solo cómo se ejecuta, el libro habrá
cumplido su propósito.

\bigskip
\begin{flushright}
{\itshape Juvenal Yosa}\\
{\sffamily\small\color{gris}2026}
\end{flushright}

\chapter*{Cómo leer este libro}
\addcontentsline{toc}{chapter}{Cómo leer este libro}
\markboth{Cómo leer este libro}{Cómo leer este libro}

Cada capítulo abre con una lista de objetivos y cierra con un resumen, una serie de ejercicios
graduados y las referencias citadas. Entre medias, el texto se apoya en un pequeño repertorio de
recuadros con funciones precisas. Conocerlos de antemano le ayudará a decidir qué leer con calma
y qué puede saltarse en una primera lectura.

\begin{intuicion}[Recuadros de intuición]
Abren cada tema con una situación cotidiana que sirve de punto de apoyo. No sustituyen la teoría:
preparan el terreno para ella.
\end{intuicion}

\begin{definicion*}{Definiciones}{}
Enuncian con precisión los objetos matemáticos o biológicos con los que trabajaremos. Están
numeradas por capítulo.
\end{definicion*}

\begin{teorema*}{Resultados}{}
Recogen las ecuaciones y los resultados centrales de cada método, normalmente seguidos de una
tabla de símbolos.
\end{teorema*}

\begin{simbolos}
\simbolo{$x$}{Cada ecuación importante va acompañada de una tabla como esta, que explica el significado de todos sus símbolos.}
\end{simbolos}

\begin{ejemplo*}{Ejemplos resueltos}{}
Aplican la teoría a números concretos, paso a paso, para que usted pueda reproducirlos con lápiz
o con Python.
\end{ejemplo*}

\begin{codigo}[ejemplo.py]
# Los fragmentos de código son breves, completos y ejecutables.
gc = sum(base in "GC" for base in "ATGCGC") / 6
\end{codigo}

\begin{historia}[Notas históricas]
Cuentan quién, cuándo y por qué se desarrolló una idea. La ciencia es una actividad humana.
\end{historia}

\begin{cuidado}[Advertencias]
Señalan errores frecuentes, supuestos que se olvidan y trampas de interpretación.
\end{cuidado}

\begin{profundizacion}[Para profundizar]
Material avanzado que puede omitirse en una primera lectura sin perder el hilo.
\end{profundizacion}

\enlacerepo{enlace al laboratorio de cada lección}

\paragraph{Convenciones} Las explicaciones están en español; el código, los nombres de funciones
y muchos términos técnicos se mantienen en inglés, como aparecen en la literatura y en la
documentación de las herramientas; su primera aparición se marca en \ingles{cursiva}. Los
decimales se escriben con coma. Las secuencias de ADN se colorean con el código de la
\cref{tab:convenciones}, el mismo que se usa en los notebooks del curso.

\begin{table}[h]
\centering\small
\caption{Colores de nucleótidos usados en todo el libro.}\label{tab:convenciones}
\begin{tabular}{@{}cccc@{}}
\toprule
\nuc{A} adenina & \nuc{C} citosina & \nuc{G} guanina & \nuc{T} timina (\nuc{U} uracilo)\\
\bottomrule
\end{tabular}
\end{table}

\paragraph{Ejercicios} Están marcados según su dificultad:
\dificultad{1}~básica, \dificultad{2}~intermedia y \dificultad{3}~avanzada. Los de nivel avanzado
suelen requerir programación o una demostración.

\paragraph{Requisitos} Se asume una familiaridad básica con cálculo, probabilidad y programación en
Python. El \cref{cap:00} repasa las herramientas computacionales y el \cref{cap:01}, la biología
molecular necesaria.

{\hypersetup{linkcolor=tinta}\tableofcontents}

\mainmatter
\setcounter{chapter}{-1}

\parte{Fundamentos}{El laboratorio, la molécula y el archivo}{%
  Antes de comparar genomas hay que saber qué es un genoma, cómo se guarda en un archivo y cómo
  se interroga con un computador. Esta parte prepara el laboratorio digital y el vocabulario
  molecular que usaremos durante todo el libro.}
\capitulo{00}
\capitulo{01}
\capitulo{02}

\parte{Comparar secuencias}{Alineamiento, perfiles y árboles}{%
  La comparación es la operación fundamental de la bioinformática. Aquí convertimos la intuición
  del parecido en algoritmos exactos, en modelos probabilísticos y, finalmente, en la historia
  evolutiva que conecta a las especies.}
\capitulo{03}
\capitulo{04}
\capitulo{05}

\parte{Del secuenciador al genoma}{Lecturas, mapeo, ensamblaje y variantes}{%
  Los secuenciadores producen millones de fragmentos imperfectos. Esta parte sigue su recorrido
  completo: la calidad de cada base, su ubicación en una referencia, la reconstrucción de genomas
  desconocidos y la detección de las diferencias que nos hacen únicos.}
\capitulo{06}
\capitulo{07}
\capitulo{08}
\capitulo{09}

\parte{Genomas en acción}{Poblaciones, transcriptomas, epigenomas y comunidades}{%
  Un genoma no es un texto estático: varía entre individuos, se expresa de forma distinta en cada
  célula, se regula químicamente y convive con miles de genomas microbianos. Aquí aprendemos a
  medir y a modelar esa dinámica.}
\capitulo{10}
\capitulo{11}
\capitulo{12}
\capitulo{13}
\capitulo{14}

\parte{Estructura, sistemas e inteligencia}{De la proteína a la red y al modelo}{%
  La última parte sube de escala: de la forma tridimensional de las proteínas a las redes que
  forman, de los modelos de aprendizaje automático a los flujos de trabajo que hacen todo esto
  reproducible.}
\capitulo{15}
\capitulo{16}
\capitulo{17}
\capitulo{18}

\backmatter
\cleardoublepage
\thispagestyle{empty}
\vspace*{\fill}
\begin{center}
\begin{tikzpicture}
  \foreach \c [count=\i] in {nucA,nucC,nucG,nucT}{\fill[\c] (\i*0.6,0) rectangle ++(0.5,0.14);}
\end{tikzpicture}\par\bigskip
{\sffamily\small\color{tinta2}
Este libro se compuso con Lua\LaTeX{} y biblatex-apa.\\
Tipografías: Libertinus Serif, Source Sans Pro y Fira Mono.\\
Figuras en TikZ y PGFPlots generadas con Python.\\[6pt]
Juvenal Yosa, PhD · Copiloto: Claude (Anthropic)\\
Licencia MIT · 2026}
\end{center}
\vspace*{\fill}

\end{document}
